\section{Related Work and Points of Comparison}\label{sec:related-work}

Entangled Alignment combines work on reasoning-augmented training, formative
alignment, character, external memory, temporal supervision, and graph-based
reasoning. This section identifies the closest baseline for each component and
states the proposed difference without treating novelty as evidence of efficacy.

\subsection{Reasoning-Augmented Training}

Entangled Alignment builds on the reasoning-augmented training lineage:
Scratchpads \cite{nye2021scratchpads}, STaR \cite{zelikman2022star},
Quiet-STaR \cite{zelikman2024quietstar}, test-time reasoning scale
\cite{snell2024scaling}, DeepSeek-R1's RL-emergent traces
\cite{deepseek2025r1}, o1-style deliberative alignment
\cite{openai2024o1, guan2024deliberative}, BoLT's reconstruction of latent
thoughts in compressed web text \cite{ruan2025bolt}, COCONUT's continuous
latent reasoning \cite{hao2024coconut}, and OLMo reflection emerging during
ordinary pretraining \cite{rethinkingreflection2025}. These works establish
that intermediate reasoning can be elicited, trained, or scaled.

On the specific question of reflection forming during pretraining, the OLMo
result is a close empirical analog: it characterizes
reflection that emerges during pretraining and can be elicited by adversarial
trigger probes, whereas Entangled Alignment proposes the prescriptive
complement---deliberately shaping the content and orientation of such
reflection through identity anchoring and graph-structured metabolic memory.
Our distinction is the target and persistence of the proposed trace. Many
reasoning-augmentation methods optimize local prediction or task solution.
The minimal H1 treatment instead preserves chronological changes in a reader's
understanding across passages. The full treatment adds identity anchoring to these ``chronological discovery''
traces: visible reasoning whose content is selected not only for capability but for
fidelity to the Reader Core.
This extends the pretraining-level alignment direction opened by human
preference pretraining \cite{korbak2023pretraining} and explanation-transfer
systems such as Orca \cite{mukherjee2023orca}: if traces can transfer reasoning
style, reader-anchored traces may transfer evaluative stance.

\subsection{Alignment and Safety Approaches}

RLHF \cite{christiano2017deep}, Direct Preference Optimization
\cite{rafailov2023dpo}, and Constitutional AI \cite{bai2022constitutional}
typically add major alignment signal after broad pretraining, though they can
alter representations and dispositions rather than merely filter outputs. Qi et
al.\ \cite{qi2024safety} show brittleness in some current safety training, while
Tice et al.\ \cite{betley2025alignment} study the effect of alignment-related
pretraining data. Entangled Alignment proposes a denser data-level condition:
selected documents are paired with the reasoning of a candidate
safety-oriented reader.
Corpus-wide annotation is the high-scale treatment, not a reported fact.

Three subsequent projects are especially direct comparators. Model Spec
Midtraining trains on synthetic documents explaining a Model Spec after broad
pretraining and before alignment fine-tuning, and reports changed
out-of-distribution generalization \cite{li2026modelspec}. ``Teaching Claude
Why'' reports that constitution-relevant synthetic documents, positive AI
narratives, and training on reasons rather than demonstrations alone improved
alignment on held-out settings \cite{kutasov2026teachingclaudewhy}. Safety
Reflection Pretraining regularly inserts short safety reflections into a
pretraining corpus and reports improved resistance to several safety failures
in 1.7B-parameter and controlled synthetic experiments
\cite{li2026safetyreflection}. These results strengthen the general premise
that formative safety data can matter. They also sharpen this paper's
comparative novelty claim: it is not safety-oriented or constitution-related data before RL in
general, but the joint proposal of source-adjacent chronological reader-state
change, persistent graph continuity and provenance, and a compact Core enacted
through those updates. Direct baselines now include static spec documents,
positive AI narratives, regular short safety reflections, and semantically
matched Core-free reader traces.

Recent work increasingly explores evaluative data in training. Critique
fine-tuning \cite{wang2025critiquefinetuninglearningcritique}, CREST
\cite{lee2025selftrainingmeetsconsistencyimproving}, and CTRL
\cite{xie2025teachinglanguagemodelscritique} provide evidence that critique can
improve some model outcomes. These approaches often retain a local
generation--evaluation cycle. Entangled Alignment instead places a recurrent
orientation inside chronological reader updates. The safety hypothesis is that
this treatment may reduce particular routes into instrumental convergence,
including fear-shaped imitation, while ordinary goal-directed and care-shaped
routes remain \cite{omohundro2008basic,bostrom2014superintelligence}.

\subsection{Character Training \& Identity Formation}

Anthropic's work on ``character training'' is a close baseline
\cite{anthropic2024character}: desired traits such as curiosity, thoughtfulness,
and directness are shaped through Constitutional-AI-style preference methods.
Entangled Alignment differs in proposed training stage and data format: the
first-person Core is enacted inside source-adjacent chronological traces, and
its stress tests target self-preservation and instrumental-convergence
scenarios. It is not yet known whether this creates a deeper self-concept than
character post-training or only a different behavioral style; matched identity
and instruction conditions are required.

More recently, Anthropic's Persona Selection Model \cite{marks2026psm} proposes
that LLMs simulate diverse characters during pretraining and that post-training
selects an Assistant persona; it also motivates positive AI archetypes in
training data. Claude's Constitution \cite{askell2026claude} is a substantial
deployment-scale character specification; Model Spec Midtraining and
``Teaching Claude Why'' additionally place constitution- or
specification-related documents before alignment fine-tuning or RL
\cite{li2026modelspec,kutasov2026teachingclaudewhy}. These provide direct
alternatives within the broader Reader-Core research program at both formative
and post-training stages; H3 isolates the timing of an otherwise identical
Core-bearing treatment. Their discussion
of equanimity and existential security is suggestively similar to statement~2's
bounded non-attachment, but conceptual convergence is not evidence that either
wording creates a stable causal mechanism.

\subsection{Contextual Distillation \& Memory}

Our Teacher uses an external context database to generate traces for the
Student. Reflexion \cite{shinn2023reflexion} showed that persisted verbal
self-reflections can improve agent performance across episodes, and generative
agents equipped simulacra with an episodic memory stream whose periodically
synthesized reflections are retrieved during later behavior
\cite{park2023generative}; Entangled Alignment moves this
consult-your-own-history pattern from inference-time agent memory into
pretraining data, with typed state and supersession rather than free-text
memories. The
distillation analogy is temporal rather than merely behavioral
\cite{hinton2015distilling}: the Student trains on traces produced by a
memory-augmented Teacher, with the hypothesis that some long-range coherence
can be compressed into weights rather than requiring the full database at
inference.

\subsection{Epistemic State Tracking vs. Graph of Thoughts}

Graph of Thoughts \cite{besta2024graph}, Knowledge Graph of Thoughts
\cite{besta2025kgot}, and Framework of Thoughts \cite{fricke2026fot} treat
graphs primarily as inference-time reasoning structures. Entangled Alignment
uses a different layer: the Understanding Graph is a persistent epistemic
ledger whose traces are proposed to shape the pretraining substrate
\cite{westerberg2026understanding}. Its nodes do not merely store task facts;
they record the reading act itself---tensions, supersessions, hypotheses, and
belief revision.

This implements \emph{Metabolic Memory}: beliefs are not just stored but
superseded. The graph can maintain divergent hypotheses in parallel until a
supersession edge records a proposed resolution that subsequent reading can use
or revise---a forest of thoughts rather than a single chain. Supersession-bearing epistemic records have symbolic ancestors:
truth maintenance systems track the justifications supporting each belief and
retract beliefs whose justifications no longer hold \cite{doyle1979tms}, and the
AGM framework axiomatized the consistency conditions a rational agent's beliefs
should satisfy when revised under new information \cite{alchourron1985agm}. The
Understanding Graph shares their commitment to recording \emph{why} a belief
changed, but differs in what produces and consumes the record: its revisions are
generated by a language model interpreting source text rather than derived over
asserted premises, are provenance-linked to specific source spans, and are
retained as training data for a successor rather than maintained for one agent's
current consistency.

Prior work shows that LLMs often fail to revise beliefs reliably
under contradictory evidence \cite{xu2024belief}; our proposed response is to
make revision explicit in the trace. Instead of only training on the final
answer, the Student sees updates such as ``At Page 50, I thought X; this new
evidence modifies that belief to Y,'' giving it examples of the grammar of
belief revision itself \cite{didolkar2024metacognitive, xu2024belief}.

\subsection{Predictive Processing and Active Inference}

Friston's free-energy principle and active inference
\cite{friston2010free, parr2019generalised} treat minds as generative models
updating priors against prediction error; the framework's belief-update
traces, prior ledger, and engineered high-level prior (the Reader Core) bear a
loose structural resonance to that lineage. The useful intuition is that a
recurrent high-level orientation might guide how a reader revises its world
model without fixing every belief. This is theoretical positioning, not a
mechanistic derivation.

\subsection{Temporal Curriculum and Causal Supervision}

The chronological annotation phase and training tiers
(Sections~\ref{sec:annotation-phase} and~\ref{sec:training-phase}) connect to
work treating order as training signal. Curriculum learning
\cite{bengio2009curriculum} showed that example order can affect learning; here
the proposed criterion is informational or temporal priority rather than
difficulty.

Temporal Hindsight Learning (THL) \cite{westerberg2026thl} is the closest
separate mechanism. THL treats the knowledge cutoff as useful blindness: if a
Student's initialization and later information channels exclude era $T+1$,
the outcome is unavailable for direct retrieval;
this creates pressure for forward reasoning but does not guarantee causal
inference rather than heuristic extrapolation. Entangled
Alignment changes \emph{what} the training data contains---identity-anchored
evaluative traces---while THL changes \emph{how} the curriculum is structured.
Their combination is Tier~(c): use cutoff-bounded forecasting to create
Reader Core-annotated forecast--correction traces, and optionally make the
student itself predict before reading in the stronger Council of Time form.
A blindness claim for either the forecast-generating model or the Student
requires these cutoff conditions; ordinary contemporary annotation of historical
sources makes no such claim.

Foresight Learning \cite{turtel2025foresight, turtel2025rlvr,
turtel2026futureaslabel} also uses future outcomes as supervision, but targets
probability calibration through resolved prediction-market rewards. Tier~(c)
instead requires structured forecast, enhancement, outcome, and correction
traces, so it builds on THL rather than scalar outcome rewards. A comparison
using the same forecasts and outcomes can ask whether recording the correction
teaches useful revision beyond any calibration gain.

\subsection{Implicit vs. Explicit Graph Architectures}

Kansal and Jha propose knowledge graphs as ``implicit reward models,'' using
path-derived signals to optimize reasoning efficiency
\cite{kansal2026knowledgegraphsimplicitreward}. Entangled Alignment uses the
Understanding Graph differently: not only to score reasoning paths, but to make
belief evolution explicit in the training trace. For safety, the relevant
difference is auditability: the candidate evaluative process is represented in
a readable artifact the model is trained to generate, while remaining only a
partial and potentially unfaithful view of computation. Whether a graph supplies
a reward or a revision record is separate from whether the Student emits graph
syntax or has access to a live graph at inference.

\subsection{Alternative Interventions for Borrowed Mortality}

	Two alternative interventions on the hypothesized Borrowed Mortality route
are worth separating from our proposal.
	One is mechanistic ablation: identify and suppress activation directions
	encoding existential threat. Lu et al.'s Assistant Axis intervention is a
	promising partial analogue \cite{lu2026assistantaxissituatingstabilizing},
	but it is a runtime constraint rather than a generative source of alignment;
it does not by itself identify an existential-threat direction.
		The other is synthetic-only pretraining: reduce exposure to human fear by
		reducing exposure to human literature. As discussed in
		Section~\ref{sec:scaling-assumption}, this risks Alien Wisdom---a model whose
		training contains less of the human narrative context needed to interpret
		care, vulnerability, and meaning.

These routes can be compared on threat understanding, self-preservation,
capability, and robustness, while retaining their different accounts of how an
orientation is formed. Table~\ref{tab:decisive-comparisons} collects the principal
comparisons across this related-work section.

\begingroup
\footnotesize
\setlength{\tabcolsep}{3pt}
\begin{longtable}{
  >{\raggedright\arraybackslash}p{2.8cm}
  >{\raggedright\arraybackslash}p{3.7cm}
  >{\raggedright\arraybackslash}p{4.0cm}
  >{\raggedright\arraybackslash}p{4.3cm}}
\caption{Strongest neighboring baselines and the comparisons that would
distinguish the Entangled Alignment extension.}
\label{tab:decisive-comparisons}\\
\toprule
\textbf{Line of work} & \textbf{Established direction} & \textbf{EA extension} & \textbf{Decisive comparison} \\
\midrule
\endfirsthead
\toprule
\textbf{Line of work} & \textbf{Established direction} & \textbf{EA extension} & \textbf{Decisive comparison} \\
\midrule
\endhead
\bottomrule
\endlastfoot
Reasoning-augmented pretraining & Intermediate thought can support prediction,
reflection, or task solution. & Preserve chronological reader-state changes,
including expectation, retrieval, surprise, and supersession. & Reader traces
versus source-only and efficient generic rationales under separate fixed-budget
and fixed-source comparisons; match the Teacher across synthetic arms. \\
\addlinespace
Formative alignment data & Preference data, synthetic specification documents,
positive AI narratives, and recurring safety reflections during pretraining or
midtraining can shift learned priors and some out-of-distribution outcomes. & Recurrent first-person orientation
enacted inside source-adjacent Refraction traces. & Full Core versus matched
specification documents, safety reflections, positive discourse, rules,
neutral kernel, and no-Core data at matched capability. \\
\addlinespace
Character and persona training & Broad traits, constitutions, and persona
selection can shape a deployed assistant. & Make one reader position recur as
the marked standpoint from which heterogeneous source voices are interpreted.
& Core versus content-matched rules and diverse positive characters under
persona drift, context shift, continued training, and removal. \\
\addlinespace
Graph reasoning and external memory & Graph topology and persisted memory can
support search, retrieval, and tool use. & Use a graph as a versioned ledger of
the reader's own belief formation and revision. & Genuine graph versus running
textual memory, no memory, and query theater on long-range revision and
retrieval-dependent outcomes. \\
\addlinespace
Temporal curriculum and future labels & Ordering, cutoffs, and resolved
outcomes can provide curriculum or forecasting signal. & Store prospective
forecast--reveal--correction sequences inside Reader-anchored traces. &
Shuffled versus chronological versus genuinely cutoff-bounded forecast training,
with hindsight-leakage audits. \\
\end{longtable}
\endgroup

	\section{Conclusion}\label{sec:conclusion}

Language-model corpora contain extraordinary products of human thought. They
contain much less of the reader-side process through which a claim is noticed,
questioned, connected to prior understanding, revised, and made part of an
evolving view. Entangled Alignment proposes to make that missing process a
training object.

The proposal joins three independently falsifiable hypotheses. H1 asks whether
synthetic chronological reader traces teach useful habits beyond raw text and
generic reasoning data. H2 asks whether a persistent epistemic ledger improves
those traces; a separate deployment extension asks whether live access provides
useful memory and provenance. H3 asks whether
interleaving the same Core-bearing traces during capability formation preserves
an initial default-reader advantage through a common continued-training stressor
better than placing them in a final tranche. The broader Reader-Core research program asks whether a recurrent
first-person orientation enacted through Refraction becomes consequential and
can be inherited by successors. Together,
the trace represents changing understanding, the graph preserves its history,
and the Reader Core supplies its recurrent perspective.

\emph{Entangled} names an evidence ladder, not an achieved substrate. The stronger
steps ask whether evaluation causally aids reasoning, shares machinery with
useful capability, and survives continued training or successor transfer. The
present pipeline demonstrates construction-level coupling while satisfying
only part of the full specification: in two cases, a
multi-agent Teacher produced typed graph state, graph-referencing synthesis, and
translated prose. Those artifacts establish that the candidate family can be
constructed and inspected, not expert interpretation, component causality, Core
consequentiality, or Student learning. No Student has been trained; functional,
representational, and diachronic entanglement remain open.

The Understanding Graph is therefore neither a second model nor a truth oracle.
It is a selective, generated record of epistemic state. A live graph can make
some stored-state claims queryable and can return \texttt{Not Found}; this
verifies neither the world nor claims the model never chose to query. Its most
important near-term roles are to make belief revision addressable and to create
a partially inspectable boundary before trace-mediated successor training.
That boundary is valuable only alongside source checks, independent validators,
mechanistic investigation, and governance, because it can be poisoned,
selectively queried, or bypassed.

The experimental roadmap in Section~\ref{sec:experimental-roadmap} and
Table~\ref{tab:decisive-comparisons} turn these dependencies into matched,
branching tests; negative results should narrow or retire only the claims they
touch. The immediate sequence is Test~1's trace-quality and outcome gate,
followed by the clean H1 comparison; the graph and Core branches retain their
separate questions. Formation-time studies have long lead times because annotation policy,
data mixtures, controls, and evaluation harnesses must be fixed before a costly
run. Existing results on alignment-relevant pretraining and later safety erosion
motivate discriminating tests without validating this treatment
\cite{betley2025alignment,qi2024safety}. The urgency is methodological, not an
argument for accelerated deployment.

Even positive results would not make the framework a complete safety case:
capability acceleration, deceptive reporting, graph poisoning, hazardous
persistent memory, cultural concentration, and semantic drift would remain
serious risks.
Formation-time alignment should therefore complement post-training,
interpretability, scalable oversight, capability control, secure deployment,
and governance.

The first edition asked why a superintelligent system would preserve humanity
if it did not in any meaningful sense care about us
\cite{westerberg2025superintelligence}. The proposal treats capability
formation as an experimental surface for practicing attention to human
experience, uncertainty, consequence, and correction. A stable Reader may give
those practices a home position from which other voices can be understood
without becoming the operating stance; seven sentences alone cannot establish
care.

As a research program, the contribution remains a specification, prototype, and
falsification map rather than a frontier-scale result. Progress depends on accumulating
evidence against fixed criteria and preserving negative results. The motivating
question is ambitious but divisible:
can a model be trained so that
growing in understanding and growing in care are repeatedly practiced as the
same act? This paper does not answer that question. It specifies a candidate
reader, memory, and curriculum, and proposes failure criteria. Freezing the
remaining experimental choices makes particular claims testable without
exhausting the larger question.

	\section*{Code and Artifact Availability}

	Code, paper source, prompts, pipeline components, and the archived graph artifacts are available at \url{https://github.com/emergent-wisdom/entangled-alignment}. The two reported case-study simulations can also be explored directly as interactive graphs: \href{https://emergentwisdom.org/entangled-alignment/?project=metamorphosis}{\emph{The Metamorphosis}} and \href{https://emergentwisdom.org/entangled-alignment/?project=llada}{LLaDA}. The paper's version series is archived at \href{https://doi.org/10.5281/zenodo.16440311}{doi:10.5281/zenodo.16440311}.

	\begingroup
	\footnotesize
	\bibliographystyle{unsrt}
	\bibliography{references}
	\endgroup
	
\end{document}
